% !TeX TS-program = lualatex
% !TeX encoding = UTF-8
\documentclass[12pt,a4paper]{article}

\usepackage{fontspec}
\setmainfont{Latin Modern Roman}
\usepackage[english]{babel}
\usepackage[a4paper,margin=25mm]{geometry}
\usepackage{booktabs}
\usepackage{tabularx}
\usepackage{hyperref}

%\setlength{\parindent}{1.2em}
%\setlength{\parskip}{0.25em}
%\newcommand{\textbf}[1]{\texttt{\detokenize{#1}}}

\day=67
\title{Race master 4D: racist game\\[0.3em]
	\large Prak documentation}
\author{Michael Gustavo \& Danila Malyshev}
\date{\today}

\begin{document}
	
	\maketitle
	
	\section{Introduction}
	
	This project is a small C++ program for viewing traffic at a four-way road intersection. It uses SFML to draw the simulation and ImGui for the control panel. The main idea is to see how traffic lights, cars, and pedestrians behave when the signal timings or traffic settings are changed.
	
	The intersection has roads in the north, south, east, and west directions. Each road has six lanes: three for incoming traffic and three for outgoing traffic. The incoming lanes have fixed movements: right turn, straight ahead, and left turn. 
	
	\section{How the simulation works}
	
	The \textbf{World} class builds the roads, traffic lights, and crosswalks. A \textbf{Road} contains six \textbf{Lane} objects. Lanes hold cars and their waypoints. \textbf{Car} and \textbf{Pedestrian} are moving objects based on the common \textbf{Entity} class. The \textbf{app} code opens the window, updates the world, and draws the objects.
	
	The intersection is in the middle of the world, and the model uses 8 pixels for one metre. A car follows the waypoints for its lane, keeps space from the car in front, and stops for a red signal. The car only reacts to the signal when it is inside the selected visibility distance. On yellow, it can continue if it is already close to the stop line. Left-turning cars can wait for opposing traffic, and cars also wait if there is no room in their exit lane. Collisions are not
	simulated.
	
	The traffic signals have phases for north--south straight/right traffic, then
	north--south left turns, followed by the same movements for east--west traffic.
	Each vehicle phase has a green and a yellow part. The last phase stops all cars
	while pedestrians use the crosswalks.
	
	There are two signal modes. In static mode, green, yellow, and the all-red
	pedestrian phase use values selected in the control panel. The default values
	are 12, 3, and 12 seconds. The pedestrian phase can last longer if somebody is
	still crossing.
	
	In automatic mode, green time depends on the queue. The target is
	\(\min(30, 5 + 1.5n)\) seconds, where \(n\) is the number of cars waiting for
	the active movement. The controller can end an empty green phase after five
	seconds if there is traffic waiting on the other road. Automatic yellow stays
	at 3 seconds. Pedestrian timing also responds to the waiting queue, with a
	target of \(\min(24, 5 + 1.5p)\) seconds, where \(p\) is the number of waiting
	pedestrians.
	
	\section{Control panel}
	
	The ImGui panel contains these settings:
	
	\begin{center}
		\begin{tabularx}{\linewidth}{@{}lX@{}}
			\toprule
			\textbf{Setting} & \textbf{What it does} \\
			\midrule
			Signal mode & Selects static or automatic timing. \\
			Green duration & Static green time: 1--90 seconds; default 12. \\
			Yellow duration & Static yellow time: 1--15 seconds; default 3. \\
			Red / walk duration & Static all-red pedestrian phase: 1--90 seconds; default
			12. \\
			Signal visibility & Distance for reacting to a signal: 5--100 metres; default
			30. \\
			Car speed range & Range for new cars: 5--160 km/h overall; default
			30--120 km/h. \\
			Car arrival interval & Minimum: 1--30 seconds; maximum: from the minimum up to
			60 seconds; default 4--8 seconds. \\
			\bottomrule
		\end{tabularx}
	\end{center}
	
	The static timing sliders only apply in static mode. The other settings can be
	changed while the simulation is running. A new car gets a speed limit sampled
	from the selected range, and starts at a speed between the minimum and its own
	speed limit.
	
	When the world starts, it tries to place one car on each road. After that, one
	random interval is selected for each cycle. At the end of the interval, the
	program tries to add one car to every road, so all roads use the same period.
	If there is no room on a road, that arrival is skipped. The \textbf{Add car}
	button tries to add one car to each incoming lane with room, and
	\textbf{Next background} changes the background image.
	
	\section{Files and running the program}
	
	The main files are:
	\begin{itemize}
		\item \textbf{main.cpp} starts the application.
		\item \textbf{app.cpp} contains the window loop, control panel, and drawing.
		\item \textbf{app.h} contains application declarations and rendering helpers.
		\item \textbf{objects.h} contains the simulation classes and most of the model
		code.
		\item \textbf{objects.cpp} includes the model header for the CMake build.
		\item \textbf{CMakeLists.txt} describes the C++ project and its dependencies.
	\end{itemize}
	
	The project uses CMake 3.24 or later, C++17, SFML 3, OpenGL, and ImGui-SFML.
	With these dependencies installed, it can be built from a terminal with:
	
	\begin{verbatim}
		cmake -S crossroads -B crossroads/build
		cmake --build crossroads/build
	\end{verbatim}
	
	The background manager looks for images under \textbf{../assets/}, but those
	images are not included in the \textbf{crossroads/} source folder. The roads and
	traffic can still be drawn by the simulation.
	
\end{document}